% ==========================
% Gói hỗ trợ màu sắc
% ==========================
\usepackage{xcolor}
\definecolor{TLgreen}{HTML}{00a79d}
\definecolor{BKblue}{HTML}{030391}
\definecolor{BKbluelight}{HTML}{1488D8}
\definecolor{modern_blue}{HTML}{007BFF}
\definecolor{modern_red}{HTML}{DC3545}
\definecolor{modern_green}{HTML}{4CAF50}
\usepackage{lastpage}
\usepackage[table,xcdraw]{xcolor}
\usepackage{tikz}
\usetikzlibrary{positioning, calc}
% \usetikzlibrary{timeline}
% ==========================
% Gói hỗ trợ toán học và định dạng
% ==========================
\usepackage{amsmath}    % Hỗ trợ các công thức toán học
\usepackage{amsfonts}   % Hỗ trợ các font toán học
\usepackage{listings}   % Hỗ trợ hiển thị mã nguồn
\lstset{
    basicstyle=\ttfamily\small,       % Font monospace nhỏ
    keywordstyle=\color{blue},        % Màu xanh cho từ khóa
    stringstyle=\color{red},          % Màu đỏ cho chuỗi
    commentstyle=\color{gray},        % Màu xám cho comment
    identifierstyle=\color{black},    % Màu đen cho tên biến/hàm
    numberstyle=\tiny\color{gray},    % Số dòng nhỏ màu xám
    backgroundcolor=\color{lightgray!20}, % Nền xám nhạt
    numbers=left,                     % Hiển thị số dòng
    frame=single,                     % Viền quanh đoạn code
    breaklines=true,                  % Tự động xuống dòng
    breakatwhitespace=true            % Chỉ xuống dòng tại khoảng trắng
}

% ==========================
% Gói hỗ trợ hình ảnh và liên kết
% ==========================
\usepackage{graphicx}                % Hỗ trợ chèn hình ảnh
\usepackage{xurl}
\usepackage[
    hidelinks,
    colorlinks=true,
    linkcolor=black,   % Màu cho link nội bộ (như mục lục)
    citecolor=black,   % Màu cho trích dẫn (nếu có)
    filecolor=black,   % Màu cho file links (ít dùng)
    urlcolor=blue      % Màu của \url
]{hyperref}
\usepackage{subcaption}  % Hỗ trợ chú thích con
\usepackage{longtable}   % Hỗ trợ bảng dài
\usepackage{multirow}
\setlength{\parindent}{0pt}          % Không thụt đầu dòng
\usepackage{parskip}                 % Thêm khoảng cách giữa các đoạn

% ==========================
% Hỗ trợ tiếng Việt và Unicode
% ==========================
\usepackage[T5]{fontenc}             % Hỗ trợ mã hóa font
\usepackage[utf8]{inputenc}          % Hỗ trợ Unicode UTF-8
\usepackage[vietnamese]{babel}          % Ngôn ngữ (vietnamese/english)
\addto\extrasvietnamese{
    \renewcommand{\equationautorefname}{Biểu thức}
    \renewcommand{\figureautorefname}{Hình}
    \renewcommand{\tableautorefname}{Bảng}
    \renewcommand{\partautorefname}{Phần}
    \renewcommand{\chapterautorefname}{Chương}
    \renewcommand{\sectionautorefname}{Mục}
    \renewcommand{\subsectionautorefname}{Mục}
    \renewcommand{\subsubsectionautorefname}{Mục}
}

% ==========================
% Thiết lập lề trang
% ==========================
\usepackage[a4paper, left=0.5in, right=1in, top=1in, bottom=1in]{geometry} % Lề trang chuẩn
\geometry{total={160mm,240mm}, top=3cm, bottom=2cm, left=2cm, right=2cm} % Ghi đè lề trang nếu cần

% ==========================
% Gói hỗ trợ bố cục nhiều cột
% ==========================
\usepackage{multicol}

% ==========================
% Gói hỗ trợ đồ họa và nền
% ==========================
\usepackage{tikz}                    % Hỗ trợ đồ họa
\usepackage{transparent}             % Hỗ trợ độ trong suốt
\usetikzlibrary{shadings, fadings}   % Thư viện TikZ cho hiệu ứng
\usetikzlibrary{calc}                % Thư viện TikZ cho tính toán tọa độ
\usetikzlibrary{positioning, fit, backgrounds, arrows.meta, shapes.geometric}
\usetikzlibrary{shapes.geometric, shapes.symbols, arrows.meta, positioning, calc, backgrounds, fit}

% ==========================
% Header và Footer
% ==========================
\usepackage{fancyhdr}                % Hỗ trợ chỉnh header và footer
\pagestyle{fancy}                    % Kích hoạt kiểu header/footer fancy
\fancyhf{}                           % Xóa header/footer mặc định
\setlength{\headheight}{15pt}

\fancypagestyle{plain}{
    \fancyhf{}
    \renewcommand{\headrulewidth}{0pt}
    \renewcommand{\footrulewidth}{1pt}
    \setlength{\headsep}{10pt}

    \fancyfoot[R]{
        {\fontsize{10pt}{15pt}\textbf{\pageFooterText}}
    }
    \setlength{\footskip}{30pt}
}

% Thiết lập header
\renewcommand{\headrulewidth}{0.4pt}

\fancyhead[L]{
    {\fontsize{10pt}{15pt}\textbf{\headerLeftText}}
}
\fancyhead[R]{
    {\fontsize{10pt}{15pt}\textbf{\supervisorName}}
}
\setlength{\headsep}{10pt}

% Thiết lập footer
\fancyfoot[R]{
    {\fontsize{10pt}{15pt}\textbf{\pageFooterText}}
}
\renewcommand{\footrulewidth}{1pt}
\setlength{\footskip}{30pt}

% Đường kẻ header/footer
% \renewcommand{\headrulewidth}{0.4pt} % Độ dày đường kẻ header
% \renewcommand{\headrule}{\color{BKblue}\hrule width\headwidth height\headrulewidth}
% \renewcommand{\footrulewidth}{0.4pt} % Độ dày đường kẻ footer

% ==========================
% Gói hỗ trợ thiết lập nền
% ==========================
\usepackage{background}
\backgroundsetup{
    scale=1,
    color=white,
    opacity=1,
    position={8cm,2.1cm},
    angle=0,
    vshift=0cm,
    pages=all,
    contents={
        \ifnum\value{page}>1
            \rule{25cm}{1cm}
        \fi
    },
}

\usepackage{float}
\usepackage{cite}
\usepackage{hyperref}

% ==========================
% Biến tùy chỉnh
% ==========================
\newcommand{\docsName}{Documentation Name}
\newcommand{\docsType}{Documentation type}
\newcommand{\authName}{
    \fontsize{20pt}{20pt}\selectfont \textit{Author 1} \\
    \fontsize{20pt}{20pt}\selectfont \textit{Author 2} \\
    \fontsize{20pt}{20pt}\selectfont \textit{Author 3} \\
}
\newcommand{\cityName}{City}
\newcommand{\orgName}{Organization}
\newcommand{\headerLeftText}{Đồ án Tốt nghiệp}
\newcommand{\supervisorName}{GVHD: Nguyễn Lý Thiên Trường}
\newcommand{\pageFooterText}{Trang \thepage/\pageref{LastPage}}

% Cho phép người dùng thay đổi giá trị các biến
\newcommand{\setDocsName}[1]{\renewcommand{\docsName}{#1}}
\newcommand{\setDocsType}[1]{\renewcommand{\docsType}{#1}}
\newcommand{\setAuthName}[1]{\renewcommand{\authName}{#1}}
\newcommand{\setCityName}[1]{\renewcommand{\cityName}{#1}}
\newcommand{\setOrgName}[1]{\renewcommand{\orgName}{#1}}
\newcommand{\setHeaderLeftText}[1]{\renewcommand{\headerLeftText}{#1}}
\newcommand{\setSupervisorName}[1]{\renewcommand{\supervisorName}{#1}}

% ==========================
% Thiết lập độ sâu của mục lục
% ==========================
\setcounter{secnumdepth}{3}    % Đánh số đến 5 cấp
\setcounter{tocdepth}{3}       % Hiển thị trong mục lục đến 3 cấp

% Định nghĩa lại tên các cấp (tùy chọn)
\renewcommand{\theparagraph}{\thesubsubsection.\arabic{paragraph}}
\renewcommand{\thesubparagraph}{\theparagraph.\arabic{subparagraph}}
